\subsection{Time integration convergence}
The scalar test is \(\dot u=-u\), \(u(0)=1\), with exact solution \(u(t)=e^{-t}\) on \(0\leq t\leq2\). Five steps from \(0.2\) through \(0.0125\) are used. The error at each sample is the exact value minus the computed value. It is measured using
\[
\|e\|_{\Delta t}
=\left(\sum_k\frac{t_{k+1}-t_k}{2}
\left(\|e_k\|_2^2+\|e_{k+1}\|_2^2\right)\right)^{1/2}.
\]
The time weights approximate an integral norm and prevent the measurement from growing merely because refinement adds samples.

A separate spherical body test sets \(J=2I_3\) and torque to zero. With constant body rate \(\omega_0\), the exact quaternion is
\[
q(t)=q(0)\otimes
\left(\cos(\|\omega_0\|t/2),
\frac{\omega_0}{\|\omega_0\|}\sin(\|\omega_0\|t/2)\right).
\]
The zero rate case is handled separately. The study uses \(\omega_0=(0.3,-0.2,0.4)^T\ \mathrm{rad/s}\), a nonidentity initial quaternion, \(t_f=4\ \mathrm{s}\), and five steps from \(0.4\) through \(0.025\ \mathrm{s}\). The attitude error is
\[
d_q=\min\bigl(\|q_{\rm exact}-q_k\|_2,\|q_{\rm exact}+q_k\|_2\bigr).
\]
This accounts for the two quaternion signs representing one orientation. Angular velocity and quaternion errors are not combined into an unscaled mixed unit norm.

Orders are estimated by fitting \(\log\|e\|_{\Delta t}=a+p\log h\). Centering \(\log h\) makes the constant and linear basis columns orthogonal under the sample inner product, so \(p\) is obtained by a projection ratio rather than a poorly conditioned polynomial fit.
\begin{table}[htbp]
\centering
\begin{tabular}{lr}
\hline
Problem and method & Fitted order\\
\hline
Scalar decay with Euler & 1.0267\\
Scalar decay with Heun & 2.0498\\
Scalar decay with RK4 & 4.0540\\
Spherical rotation with RK4 & 3.9995\\
\hline
\end{tabular}

\caption{Measured orders from the saved five grid studies. Finite step effects account for small departures from the asymptotic integer orders.}
\end{table}
\begin{figure}[htbp]
\centering
\includegraphics[width=0.82\linewidth]{verification_convergence.png}
\caption{Exact scalar decay convergence. Euler, Heun, and RK4 approach first, second, and fourth order respectively. The vertical axis uses the time weighted error norm.}
\end{figure}

\subsection{Conservation over finite trajectories}
The asymmetric torque free study uses \(J=\operatorname{diag}(0.030,0.035,0.007)\ \mathrm{kg\,m^2}\), initial rate \((0.4,0.3,0.5)^T\ \mathrm{rad/s}\), and a duration of \(20\ \mathrm{s}\). Energy and the full inertial momentum vector are compared with their initial values. These are independent physical invariants rather than a finer numerical trajectory used as truth.
\begin{figure}[htbp]
\centering
\includegraphics[width=0.82\linewidth]{verification_invariants.png}
\caption{Maximum relative energy and inertial momentum errors for the asymmetric torque free body. Both decrease under refinement. Their measured powers need not coincide on every grid.}
\end{figure}
Quaternion correction norms are saved separately. Projection controls the unit constraint but does not by itself establish correct rotational dynamics.

\subsection{Linear stability checks}
For \(\lambda=-1\), the implemented RK4 step agrees with its analytic amplification polynomial. In particular, \(R(-2)=1/3\) gives decay, whereas \(R(-3)=1.375\) gives numerical growth for the same decaying differential equation. This checks stability separately from convergence order.
\begin{figure}[htbp]
\centering
\includegraphics[width=0.82\linewidth]{verification_stability.png}
\caption{Amplification magnitude on the negative real axis. The horizontal line is one. Values below it represent decay for the scalar linear model. A stable factor alone does not establish accuracy.}
\end{figure}

\subsection{Scope of special-case verification}
These are exact solution and conservation studies for stated special cases.
The following subsections document the real field and controller checks,
saturation-aware refinement, and the accepted $0.25\ \mathrm{s}$ timestep.
No engineering success fraction is inferred from these special cases;
the engineering results use the separate 16-case real targeted study.
